\subsubsection{Linear Discriminant Analysis}
Linear Discriminant Analysis (LDA) có thể được hiểu dưới hai góc nhìn tương
đương: góc nhìn xác suất (probabilistic) và góc nhìn hình học (projection).

\paragraph{Góc nhìn xác suất}
LDA giả định rằng dữ liệu của mỗi lớp tuân theo phân phối Gaussian đa biến với
cùng ma trận hiệp phương sai:

\begin{equation}
  \boldsymbol{x} \mid t = k \sim \mathcal{N}(\boldsymbol{\mu}_k,
  \boldsymbol{\Sigma}).
\end{equation}
Dựa trên định lý Bayes, ta thu được hàm phân biệt tuyến tính:

\begin{equation}
  \delta_k(\boldsymbol{x}) = \boldsymbol{x}^\intercal \boldsymbol{\Sigma}^{-1}
  \boldsymbol{\mu}_k - \frac{1}{2} \boldsymbol{\mu}_k^\intercal
  \boldsymbol{\Sigma}^{-1} \boldsymbol{\mu}_k + \log \pi_k.
\end{equation}
Mẫu $\boldsymbol{x}$ được gán vào lớp có giá trị $\delta_k(\boldsymbol{x})$
lớn nhất. Do các lớp chia sẻ cùng $\boldsymbol{\Sigma}$, biên quyết định là
tuyến tính.

\paragraph{Góc nhìn hình học}

Từ góc nhìn hình học, LDA tìm một phép chiếu tuyến tính
$\boldsymbol{w}$ sao cho khi chiếu dữ liệu:
\begin{equation}
  z = \boldsymbol{w}^\intercal \boldsymbol{x},
\end{equation}
các lớp được phân tách tốt nhất.

Tiêu chí tối ưu là cực đại hóa tỷ lệ giữa phương sai giữa các lớp và
phương sai trong lớp:
\begin{equation}
  \boldsymbol{w}^* = \arg\max_{\boldsymbol{w}}
  \frac{\boldsymbol{w}^\intercal \boldsymbol{S}_B \boldsymbol{w}}
       {\boldsymbol{w}^\intercal \boldsymbol{S}_W \boldsymbol{w}}.
\end{equation}

Trong đó:

\begin{itemize}
  \item $\boldsymbol{S}_B$ (between-class variance) đo mức độ phân tách
  giữa các lớp, được định nghĩa:
  \begin{equation}
    \boldsymbol{S}_B =
    \sum_{k=1}^{K} N_k
    (\boldsymbol{\mu}_k - \boldsymbol{\mu})
    (\boldsymbol{\mu}_k - \boldsymbol{\mu})^\intercal,
  \end{equation}
  với:
  \begin{itemize}
    \item $K$ là số lớp
    \item $N_k$ là số mẫu của lớp $k$
    \item $\boldsymbol{\mu}_k$ là trung bình của lớp $k$
    \item $\boldsymbol{\mu}$ là trung bình toàn bộ dữ liệu
  \end{itemize}

  \item $\boldsymbol{S}_W$ (within-class variance) đo mức độ phân tán
  của các điểm trong từng lớp:
  \begin{equation}
    \boldsymbol{S}_W =
    \sum_{k=1}^{K} \sum_{\boldsymbol{x}_i \in C_k}
    (\boldsymbol{x}_i - \boldsymbol{\mu}_k)
    (\boldsymbol{x}_i - \boldsymbol{\mu}_k)^\intercal,
  \end{equation}
  trong đó $C_k$ là tập các mẫu thuộc lớp $k$.
\end{itemize}

Trực quan:
\begin{itemize}
  \item $\boldsymbol{S}_B$ càng lớn $\Rightarrow$ các lớp càng tách xa nhau
  \item $\boldsymbol{S}_W$ càng nhỏ $\Rightarrow$ các điểm trong lớp càng
  tập trung
\end{itemize}

Do đó, LDA tìm $\boldsymbol{w}$ sao cho:
\begin{itemize}
  \item khoảng cách giữa các lớp là lớn nhất
  \item độ phân tán trong mỗi lớp là nhỏ nhất
\end{itemize}

\paragraph{Nhận xét}

Hai cách tiếp cận trên là tương đương và dẫn đến cùng một nghiệm. Góc nhìn xác
suất giúp diễn giải LDA như một mô hình sinh (generative model), trong khi góc
nhìn hình học cho thấy LDA là một phương pháp giảm chiều có giám sát, tối đa
hóa khả năng phân tách giữa các lớp.

Trong thực tế, LDA hoạt động tốt khi dữ liệu gần với giả định Gaussian và các
lớp có phương sai tương tự nhau, nhưng có thể kém hiệu quả khi các giả định
này bị vi phạm hoặc khi quan hệ giữa các lớp là phi tuyến.
